\subsection{\texorpdfstring{$\mathcal M_2$ $3$-水平边界上的共同 $24$ 维可见核与 $15$ 维缺陷块}{M2 3-水平边界上的共同 24 维可见核与 15 维缺陷块}}

正文关于 genus-\(2\) level-\(3\) Klingen/Siegel/旗标覆盖的边界分析，已经分别给出点--Lagrangian 入射对应所诱导的共同 \(24\) 维块、\(\Delta_0\) 与 \(\Xi\) 处的惰性特征多项式，以及旗标覆盖在 \(\Delta_0\) 上对消失向量方向的精细分裂。
把这些结果重新并置后，可把两类 \(40\) 重覆盖的全部差异压缩为一条“共同 \(24\) 维核 + 纯 \(15\) 维缺陷”主线。

\begin{theorem}[共同 \(24\) 维可见核与纯 \(15\) 维缺陷差]\label{thm:derived-m2-level3-common24-defect15-exact-sequence}
\leanverified{paper\_derived\_m2\_level3\_common24\_defect15\_exact\_sequence}
记两类秩 \(40\) 置换局部系统为
\[
\mathbb V_{\mathrm{Kl}},
\qquad
\mathbb V_{\mathrm{Si}},
\]
以及它们的 Hecke 本征分解
\[
\mathbb V_{\mathrm{Kl}}=\QQ\oplus \mathbb V_{24}\oplus \mathbb V_{15}^{\mathrm{Kl}},
\qquad
\mathbb V_{\mathrm{Si}}=\QQ\oplus \mathbb V_{24}\oplus \mathbb V_{15}^{\mathrm{Si}}.
\]
则存在自然短正合列
\[
0\longrightarrow \mathbb V_{15}^{\mathrm{Kl}}
\longrightarrow \mathbb V_{\mathrm{Kl}}
\xrightarrow{\ \Phi\ }
\mathbb V_{\mathrm{Si}}
\longrightarrow \mathbb V_{15}^{\mathrm{Si}}
\longrightarrow 0,
\]
并且
\[
\mathrm{Im}(\Phi)=\QQ\oplus \mathbb V_{24}.
\]
因此，两类 \(40\) 重覆盖的共同可见内容被唯一压缩到 \(\QQ\oplus \mathbb V_{24}\) 中，而其全部差异严格集中在两条互不同构的 \(15\) 维缺陷块上。

在 \(\Delta_0\) 的一般点处，惰性生成元 \(\tau\) 满足
\[
\chi_{\tau|\mathbb V_{24}}(X)=(X-1)^{12}(X^2+X+1)^6
\]
于 Klingen 与 Siegel 两侧完全相同，而
\[
\chi_{\tau|\mathbb V_{15}^{\mathrm{Kl}}}(X)=(X-1)^9(X^2+X+1)^3,
\]
\[
\chi_{\tau|\mathbb V_{15}^{\mathrm{Si}}}(X)=(X-1)^3(X^2+X+1)^6.
\]
特别地，
\[
\Tr(\tau|\mathbb V_{15}^{\mathrm{Kl}})=6,
\qquad
\Tr(\tau|\mathbb V_{15}^{\mathrm{Si}})=-3,
\]
且驯导子指数满足
\[
a_{\Delta_0}(\mathbb V_{15}^{\mathrm{Kl}})=6,
\qquad
a_{\Delta_0}(\mathbb V_{15}^{\mathrm{Si}})=12.
\]
\end{theorem}

\begin{proof}
正合列与像的描述由定理
\ref{thm:conclusion-m2-level3-incidence-24-identification-kill-minus4}
及推论 \ref{cor:conclusion-m2-level3-incidence-exact-sequence-15-separation}
直接给出。
\(\Delta_0\) 处 \(\mathbb V_{24}\) 与两条 \(15\) 维块的惰性特征多项式分别来自定理
\ref{thm:conclusion-m2-level3-delta0-inertia-klingen-charpoly}
与定理 \ref{thm:conclusion-m2-level3-delta0-inertia-siegel-charpoly}；
导子指数则由推论
\ref{cor:conclusion-m2-level3-delta0-inertia-invariants-conductor-determinant}
读出。证毕。
\end{proof}

\leanverified{paper\_derived\_m2\_level3\_subgroup\_index\_controls\_degeneration}
\begin{corollary}[转射碰撞边界的 Jacobian 退化由子群指数控制]\label{cor:derived-m2-level3-subgroup-index-controls-degeneration}
设 \(\Delta_0\) 边界的一参数退化由两条相邻换位局部单值 \(g_i,g_{i+1}\in S_4\) 的碰撞给出，并记
\[
H:=\langle g_i,g_{i+1}\rangle\in\{C_2,V_4,S_3\},
\qquad
\delta:=[S_4:H]\in\{12,6,4\}.
\]
则对应半稳定中心纤维由一个主分量与 \(\delta\) 个有理尾分量组成；其对偶图具有 \(\delta\) 个独立圈，从而全部节点都贡献非分离退化。
主分量 \(X_{\mathrm{main}}\) 的亏格满足
\[
g(X_{\mathrm{main}})=16-\delta\in\{4,10,12\},
\]
并且极限 Jacobian 处于半阿贝尔精确列
\[
0\longrightarrow \mathbf G_m^\delta
\longrightarrow J_{\lim}
\longrightarrow J(X_{\mathrm{main}})
\longrightarrow 0.
\]
因此，三类边界中的 \(S_3\)-型同时具有最小的环面秩与最大的阿贝尔主分量，而 \(C_2\)-型则在相反方向达到极值。
\end{corollary}

\begin{proof}
由定理 \ref{thm:xi-s4-boundary-vanishing-cycle-induced-module}，
消失圈 \(S_4\)-模同构于 \(\Ind_H^{S_4}(\xi)\)，其维数恰为 \([S_4:H]=\delta\)；同一定理的证明又给出中心纤维由一个主分量与 \(\delta\) 个有理尾分量组成，且每个尾分量与主分量形成一条图圈，故节点皆贡献非分离退化，极限 Jacobian 的环面秩即为 \(\delta\)。
再由推论 \ref{cor:xi-s4-boundary-main-component-h10-decomposition}，
三类边界上 \(H^{1,0}(X_{\mathrm{main}})\) 的维数分别为 \(4,10,12\)，亦即
\(
g(X_{\mathrm{main}})=16-\delta
\)。
最后对半稳定曲线应用广义 Jacobian 的标准结构定理，得到所示半阿贝尔扩张。证毕。
\end{proof}

\begin{theorem}[旗标覆盖在 \texorpdfstring{$\Delta_0$}{Delta0} 上唯一记住消失向量]\label{thm:derived-m2-level3-delta0-flag-vanishing-line-memory}
\leanverified{paper\_derived\_m2\_level3\_delta0\_flag\_vanishing\_line\_memory}
记遗忘态射
\[
p:\overline{\mathcal M}_2^{\mathrm{Fl}}(3)\longrightarrow \overline{\mathcal M}_2^{\mathrm{Kl}}(3).
\]
在 \(\Delta_0\) 的一般点处，记 Klingen 侧分解为
\[
\bigl(\overline{\mathcal M}_2^{\mathrm{Kl}}(3)\bigr)^{-1}(\Delta_0)
=
\Bigl(\bigsqcup_{i=1}^{13}D_i^{(1)}\Bigr)\ \sqcup\ \Bigl(\bigsqcup_{j=1}^{9}D_j^{(3)}\Bigr).
\]
则存在且仅存在一个特异的无分歧分量
\[
D_\star^{(1)},
\]
它对应于转射 \(\tau_v\) 的消失向量所张成的射影点 \([\langle v\rangle]\in \PP(v^\perp)\)，并满足
\[
p^{-1}(D_\star^{(1)})=\bigsqcup_{a=1}^{4}F_{\star,a}^{(1)},
\]
其中四个 lifting 在一般点处全部相对无分歧。

对其余每个无分歧分量 \(D_i^{(1)}\neq D_\star^{(1)}\)，都有严格分裂
\[
p^{-1}(D_i^{(1)})=F_i^{(1)}\sqcup F_i^{(3)},
\]
其中一支相对无分歧，另一支在一般点处相对三次分歧。
而对每个三次分歧分量 \(D_j^{(3)}\)，有
\[
p^{-1}(D_j^{(3)})=\bigsqcup_{b=1}^{4}F_{j,b}^{(3)},
\]
且这些 lifting 相对 \(D_j^{(3)}\) 不再新增分歧。
因此，旗标覆盖并非仅作次数 \(4\) 的细化，而是精确隔离出唯一的消失向量方向。
\end{theorem}

\begin{proof}
这正是定理 \ref{thm:conclusion-m2-level3-flag-over-klingen-delta0-refined-splitting}
的内容。
在 \(\Delta_0\) 一般点处，惯性生成元可取为转射 \(\tau_v\)；若 Klingen 直线正好等于 \(\langle v\rangle\)，则其上方四个旗标全部固定，因而给出唯一的“四支全无分歧”情形。
其余无分歧直线与三次分歧直线则分别产生
\(
1+3
\)
与
\(
3+3+3+3
\)
型 lifting 分裂。证毕。
\end{proof}

\begin{theorem}[共同 \(24\) 维块的跨边界稳定性与 \(15\) 维缺陷块的边界敏感性]\label{thm:derived-m2-level3-common24-crossboundary-stability}
\leanverified{paper\_derived\_m2\_level3\_common24\_crossboundary\_stability}
沿定理 \ref{thm:derived-m2-level3-common24-defect15-exact-sequence} 的记号。
则共同 \(24\) 维块 \(\mathbb V_{24}\) 具有跨边界稳定性：在 \(\Delta_0\) 的一般点处，
\[
\chi_{\tau|\mathbb V_{24}}(X)=(X-1)^{12}(X^2+X+1)^6
\]
于 Klingen/Siegel 两侧完全相同；而在双椭圆除子 \(\Xi\) 的一般点处，惯性生成元 \(\sigma\) 满足
\[
\chi_{\sigma|\mathbb V_{24}}(T)=(T-1)^{16}(T+1)^8,
\]
仍然在两侧完全一致。

相反，两条 \(15\) 维缺陷块在 \(\Xi\) 上继续分叉：
\[
\chi_{\sigma|\mathbb V_{15}^{\mathrm{Kl}}}(T)=(T-1)^7(T+1)^8,
\]
\[
\chi_{\sigma|\mathbb V_{15}^{\mathrm{Si}}}(T)=(T-1)^{11}(T+1)^4.
\]
因此，\(\mathbb V_{24}\) 是同时对转射型边界与对合型边界稳定的最大可见核心，而一切边界敏感信息都被压缩在两条 \(15\) 维缺陷块中。
\end{theorem}

\begin{proof}
\(\Delta_0\) 处的共同特征多项式已经在定理
\ref{thm:derived-m2-level3-common24-defect15-exact-sequence}
中给出。
\(\Xi\) 处的共同 \(24\) 维块与两条 \(15\) 维块的惰性特征多项式则由定理
\ref{thm:conclusion-m2-level3-xi-inertia-hecke-eigensystems-charpoly}
直接给出。
将两类边界并读即可看出：\(\mathbb V_{24}\) 在两侧始终保持一致，而所有差异都继续集中在
\(\mathbb V_{15}^{\mathrm{Kl}}\) 与 \(\mathbb V_{15}^{\mathrm{Si}}\) 上。证毕。
\end{proof}

\begin{remark}[共同核与缺陷块的读法]\label{rem:derived-m2-level3-common24-defect15-reading}
上述三条结果共同说明：Klingen 与 Siegel 两类 \(40\) 重覆盖之间并不存在“处处微妙不同”的全局分叉。
真正稳定的可见部分被一次性压缩到 \(\mathbb V_{24}\)，而所有边界敏感内容都纯化为两条 \(15\) 维缺陷块，并可在 \(\Delta_0\) 与 \(\Xi\) 两类边界上被惰性谱分别分辨。
\end{remark}

\endinput
